\documentclass[11pt]{article}

% Page layout
\usepackage[a4paper,margin=1in]{geometry}
\usepackage{hyperref}

% Fonts and formatting
\usepackage{times}
\usepackage{setspace}
\usepackage{titlesec}

% Math, figures, tables
\usepackage{amsmath,amssymb}
\usepackage{graphicx}
\usepackage{booktabs}
\usepackage{array}
\usepackage{multirow}

% Code listings
\usepackage{listings}
\usepackage{xcolor}
\lstset{
    basicstyle=\ttfamily\footnotesize,
    breaklines=true,
    columns=fullflexible,
    keywordstyle=\color{blue!70!black},
    commentstyle=\color{gray},
    stringstyle=\color{red!60!black},
    showstringspaces=false,
    frame=single,
    framesep=4pt,
}

% Citations and links
\usepackage{cite}

% Section style
\titleformat{\section}{\large\bfseries}{\thesection}{1em}{}
\titleformat{\subsection}{\normalsize\bfseries}{\thesubsection}{1em}{}

% Line spacing
\setstretch{1.1}

\title{\textbf{Graphormer: Inference-Time Optimisation Ablation \\
\large CS768 Course Project --- End-Semester Report}}

\author{
Harshil Amitkumar Solanki (23B1016), Sumedh S S (23B1079) \\
}

\date{}

\begin{document}

\maketitle

%-----------------------------------------------------------------------------
\begin{abstract}
Following our mid-semester study of the Graphormer architecture
\cite{graphormer}, we now report an empirical ablation of inference-time
optimisations applied to the model. Because the official Microsoft codebase
is pinned to PyTorch~1.9.1 / Python~3.9 / Fairseq, and predates fused
attention kernels, we re-implemented the Graphormer encoder, structural
encodings (centrality, spatial, edge), and graph-prediction head as a
standalone pure-PyTorch~2.x module with a Fairseq-compatible checkpoint
loader. On top of this we benchmarked three attention back-ends ---
hand-written \texttt{matmul+softmax}, PyTorch's \emph{scaled\_dot\_product\_attention}
(SDPA), and explicit FlashAttention-2 dispatch --- across three precisions
(fp32, fp16, bf16) on an RTX~3050 (Ampere) GPU. For the
\texttt{graphormer\_base} configuration ($\sim$47M params) the optimised path
yields a $\mathbf{2.09\times}$ end-to-end speed-up over the eager baseline,
while for the \texttt{graphormer\_slim} configuration ($\sim$0.96M params) the
same optimisations are net-negative because kernel-launch overhead dominates.
We also analyse why KV-caching --- a standard decoder-side optimisation ---
is mathematically and structurally inapplicable to a Graphormer-style
encoder, and exclude it with justification. Code:
\url{https://github.com/HARSHIL2836R/Graphormer-Project.git}.
\end{abstract}

%-----------------------------------------------------------------------------
\section{Introduction}

Graphs are a fundamental data structure for representing relational
information across domains such as molecular chemistry, citation networks
and social systems. Graph Neural Networks (GNNs) are the dominant family
of learning models for these data, but most GNN variants assume fixed
structures and homogeneous node/edge types. Graphormer~\cite{graphormer}
adapts the Transformer architecture to graphs by injecting graph topology
into self-attention via three structural encodings: centrality, spatial
position, and edge features.

Our mid-sem report covered the formulation, the structural encodings, and
related work (GTN~\cite{gtn}, PAGTN~\cite{pagtn}, GRPE~\cite{grpe},
SAN~\cite{san}). For the end-sem deliverable we shift focus from
\emph{what} the model computes to \emph{how efficiently} it can be run at
inference time. Concretely, we ask: when modern attention back-ends
(SDPA, FlashAttention-2) and reduced precisions (fp16, bf16) are bolted
onto Graphormer's structurally-biased attention, how much latency and VRAM
do we actually save, and where does the win come from? The result is
non-trivial because Graphormer's attention is not vanilla: every layer
adds a learned bias matrix
$B \in \mathbb{R}^{H \times (N{+}1) \times (N{+}1)}$ inside the softmax,
which constrains which fused kernels are eligible.

%-----------------------------------------------------------------------------
\section{Problem Formulation}

Let $G = (V,E)$ be a graph with node features $X \in \mathbb{R}^{|V| \times d}$
and adjacency $A \in \{0,1\}^{|V| \times |V|}$. Graphormer encodes the
graph as a sequence $[\textsc{cls}, v_1, \dots, v_N]$ and applies $L$
Transformer encoder layers with a structurally-biased attention,
\begin{equation}
\mathrm{Attn}(Q,K,V) = \mathrm{softmax}\!\left(\frac{QK^{\top}}{\sqrt{d_k}} + B\right) V,
\end{equation}
where $B$ aggregates centrality, shortest-path spatial encoding and edge
features. For graph-level regression / classification, the [CLS] token's
final representation is passed through an output head.

The end-sem study fixes the model and asks the following \emph{inference}
question. Let $\mathcal{F}$ denote the forward pass and let
$\theta \in \{\textit{backend}, \textit{precision}\}$ be the optimisation
configuration. Define
\[
\mathcal{L}(\theta) = \mathbb{E}\big[\mathrm{latency}(\mathcal{F}_\theta)\big],
\qquad
\mathcal{M}(\theta) = \max\big[\mathrm{VRAM}(\mathcal{F}_\theta)\big].
\]
We want
$\arg\min_{\theta} \mathcal{L}(\theta)$ and the corresponding
$\mathcal{M}(\theta)$, subject to the constraint that the structural bias
$B$ remain numerically intact (i.e.\ optimisations that drop or distort the
bias are disallowed).

%-----------------------------------------------------------------------------
\section{Related Work (recap)}

We retain the four works surveyed at mid-sem: GTN \cite{gtn} (learned
meta-paths), PAGTN \cite{pagtn} (path-augmented attention), GRPE
\cite{grpe} (relative positional encoding for graph Transformers) and
SAN \cite{san} (Laplacian-based structural positional encoding). All four
support the central claim of Graphormer that Transformers can match or
exceed GNNs once a sufficient structural prior is injected into the
attention. The end-sem extension is orthogonal: rather than proposing a
new structural encoding, we evaluate how Graphormer's existing encoding
interacts with modern attention kernels.

%-----------------------------------------------------------------------------
\section{Methodology}

\subsection{Standalone Re-implementation}

The official Microsoft Graphormer~\cite{graphormer-msft} ships as a
Fairseq~\cite{fairseq} extension pinned to PyTorch~1.9.1 + CUDA~11.1 +
Python~$\le$3.9, with transitive dependencies on \texttt{dgl}~0.7.2,
\texttt{torch-scatter}/\texttt{torch-sparse}, \texttt{ogb}~1.3.2 and
\texttt{rdkit-pypi}~2021.9.3. The host machine for this study runs
Ubuntu~24.04 with Python~3.12 and a CUDA-13 driver, and PyTorch~1.9.1
predates the introduction of
\texttt{torch.nn.functional.scaled\_dot\_product\_attention} (added in
PyTorch~2.0). A native install was therefore not viable, and even if it
were, it could not host the kernels under study.

We re-implemented the model in pure PyTorch~$\ge$2.1, preserving the
module hierarchy and tensor shapes of the original code so that the
official checkpoints transfer without re-training. The five extracted
components are:
\begin{itemize}
    \item \textbf{GraphNodeFeature} --- atom-feature embedding (sum over
        feature slots) plus in/out-degree centrality embeddings, prefixed
        with a learnable [CLS] token.
    \item \textbf{GraphAttnBias} --- spatial-position embedding,
        multi-hop edge-feature aggregation along shortest paths, and
        the learned [CLS]-virtual-distance bias, all combined into the
        per-head additive bias $B$.
    \item \textbf{MultiheadAttention} --- with a switchable back-end
        (Sec.~\ref{sec:backends}).
    \item \textbf{GraphormerEncoderLayer} --- post- or
        pre-LayerNorm transformer block (post-LN by default, matching
        \texttt{graphormer\_base}).
    \item \textbf{GraphormerForGraphPrediction} --- the [CLS] read-out
        head matching the Fairseq \texttt{GraphormerEncoder}.
\end{itemize}
A checkpoint loader strips the Fairseq \texttt{encoder.}~prefix and
splits the legacy \texttt{in\_proj\_\{weight,bias\}} buffers (used by
old PyTorch \texttt{nn.MultiheadAttention}) into separate
$Q,K,V$ projections, so the official
\texttt{checkpoint\_best\_pcqm4mv1.pt} drops in without manual surgery.

\subsection{Attention Back-ends Studied}
\label{sec:backends}

The forward pass of self-attention with structural bias is implemented
three different ways while the input/output tensor shapes are held
identical:

\paragraph{(1) \texttt{manual}.} Explicit \texttt{torch.matmul} +
\texttt{softmax}, mirroring the original \texttt{MultiheadAttention} from
\texttt{graphormer/modules/multihead\_attention.py}. The query is
pre-scaled by $1/\sqrt{d_k}$, padding is folded into $B$ as $-\infty$
columns, and the dot-product score is materialised as a
$B{\times}H{\times}T{\times}T$ tensor before softmax. This is the
numerically faithful baseline.

\paragraph{(2) \texttt{sdpa}.} A single call to
\texttt{torch.nn.functional.scaled\_dot\_product\_attention} with the
combined bias passed as \texttt{attn\_mask}. The scaling is applied
internally, so the explicit \texttt{$\times 1/\sqrt{d_k}$} step is
removed. PyTorch dispatches at runtime to the fastest eligible kernel
(FlashAttention-2 / xFormers memory-efficient / math fall-back), which
on Ampere with an additive float bias is typically the
memory-efficient kernel.

\paragraph{(3) \texttt{sdpa-flash}.} Forces FlashAttention-2 as the
preferred back-end via the \texttt{torch.nn.attention.sdpa\_kernel}
context manager, with memory-efficient as the fallback. FlashAttention-2
admits an additive bias only on a subset of (head\_dim, seq\_len, dtype)
configurations; the context-manager fallback path catches the remaining
cases.

For all three back-ends, the structural bias $B$ is identical and is
\emph{not} approximated, ensuring that any latency/memory difference is
attributable purely to the kernel and not to the model.

\subsection{Why KV-caching is excluded}
\label{sec:nokv}

KV-caching accelerates \emph{autoregressive} decoders by re-using the
cached $K, V$ projections of all previously-emitted tokens when the
$(t{+}1)$-th token is generated, so each new step costs $O(t)$ instead of
$O(t^2)$. Two structural facts make this inapplicable to Graphormer:
\begin{enumerate}
    \item \textbf{Encoder-only, single bidirectional pass.} The model
        consumes the full padded graph $[\textsc{cls}, v_1, \dots, v_N]$
        in one forward, with no autoregressive step over which $K,V$
        could be reused. The [CLS] read-out is a single shot.
    \item \textbf{Globally coupled bias.} The structural bias $B$ is a
        function of all node-pairs jointly --- spatial encoding requires
        the full shortest-path matrix, edge encoding requires the
        multi-hop tensor along those paths. There is no time-axis along
        which $B$ factorises, so even a hypothetical incremental scheme
        could not partially reuse prior $K,V$.
\end{enumerate}
We therefore exclude KV-caching from the ablation and document it
explicitly in the codebase rather than reporting a misleading null
result.

\subsection{Experimental Setup}

\paragraph{Hardware.} NVIDIA RTX~3050 Laptop GPU (Ampere, sm\_86), 4~GB
VRAM, on a single workstation running Ubuntu~24.04 with NVIDIA
driver~580 / CUDA-13 runtime exposure.

\paragraph{Software.} Python~3.11, PyTorch~2.5.1 + CUDA-12.1 wheels
(forward-compatible with the cu13 driver), no \texttt{torch.compile}, no
graph-level fusion. All runs are single-process, single-GPU, no AMP
gradient scaling (inference-only).

\paragraph{Workload.} Synthetic mini-batches matching the schema
produced by the Graphormer collator: \texttt{x}~$\in\!\mathbb{Z}^{B \times N \times K}$
atom-feature ids, \texttt{in/out\_degree}~$\in\!\mathbb{Z}^{B \times N}$,
\texttt{spatial\_pos}~$\in\!\mathbb{Z}^{B \times N \times N}$,
\texttt{attn\_edge\_type}~$\in\!\mathbb{Z}^{B \times N \times N \times K_e}$,
\texttt{edge\_input}~$\in\!\mathbb{Z}^{B \times N \times N \times D_{\text{max}} \times K_e}$.
All ids are sampled uniformly within the embedding-table ranges, so
embedding lookups never go out of bounds and computational cost is
identical to a real ZINC / OGBG-MolHIV batch of the same shape.

\paragraph{Configurations.} Two architectures from the original code:
\texttt{graphormer\_slim} (12L, $d{=}80$, $h{=}8$, $\sim$0.96M params)
benchmarked at $B{=}8, N{=}32, 50$ timed iterations after $10$ warm-up
iterations; and \texttt{graphormer\_base} (12L, $d{=}768$, $h{=}32$,
$\sim$47M params) benchmarked at $B{=}2, N{=}32, 30$ timed iterations
after $5$ warm-up. Smaller $B$ for the base model is required by the
4~GB VRAM budget once fp16/bf16 autocast doubles activation memory.

\paragraph{Metrics.} Wall-clock latency of one forward pass under
\texttt{torch.no\_grad()}, measured with
\texttt{torch.cuda.synchronize()} on either side of
\texttt{time.perf\_counter}; peak GPU memory via
\texttt{torch.cuda.max\_memory\_allocated} after a
\texttt{reset\_peak\_memory\_stats}. We report mean, median, p90, and
standard deviation of latency, plus peak-memory in MiB.

%-----------------------------------------------------------------------------
\section{Results}

\subsection{\texttt{graphormer\_base} (47M params, $B{=}2$, $N{=}32$)}

\begin{table}[h!]
\centering
\small
\begin{tabular}{llrrrrrr}
\toprule
\textbf{Backend} & \textbf{Dtype} & \textbf{Mean (ms)} & \textbf{Median} & \textbf{p90} &
\textbf{Std} & \textbf{Peak (MiB)} & \textbf{$\times$ baseline} \\
\midrule
manual     & fp32 & 5.376 & 5.371 & 5.391 & 0.015 & 198.79 & 1.00$\times$ \\
manual     & fp16 & 2.910 & 2.846 & 2.920 & 0.293 & 292.52 & 1.85$\times$ \\
manual     & bf16 & 2.893 & 2.817 & 2.856 & 0.329 & 292.52 & 1.86$\times$ \\
sdpa       & fp32 & 4.771 & 4.776 & 4.823 & 0.056 & 200.04 & 1.13$\times$ \\
sdpa       & fp16 & \textbf{2.567} & \textbf{2.442} & \textbf{2.491} & 0.480 & 292.52 & \textbf{2.09$\times$} \\
sdpa       & bf16 & 2.593 & 2.527 & 2.573 & 0.329 & 292.52 & 2.07$\times$ \\
sdpa-flash & fp32 & 4.714 & 4.701 & 4.826 & 0.081 & 200.04 & 1.14$\times$ \\
sdpa-flash & fp16 & 2.641 & 2.580 & 2.636 & 0.294 & 292.52 & 2.04$\times$ \\
sdpa-flash & bf16 & 2.670 & 2.595 & 2.716 & 0.330 & 292.52 & 2.01$\times$ \\
\bottomrule
\end{tabular}
\caption{Inference latency and peak VRAM for \texttt{graphormer\_base}.
Baseline (1.00$\times$) is \texttt{manual / fp32}.}
\label{tab:base}
\end{table}

\subsection{\texttt{graphormer\_slim} (0.96M params, $B{=}8$, $N{=}32$)}

\begin{table}[h!]
\centering
\small
\begin{tabular}{llrrrrrr}
\toprule
\textbf{Backend} & \textbf{Dtype} & \textbf{Mean (ms)} & \textbf{Median} & \textbf{p90} &
\textbf{Std} & \textbf{Peak (MiB)} & \textbf{$\times$ baseline} \\
\midrule
manual     & fp32 & \textbf{2.513} & \textbf{2.499} & \textbf{2.561} & 0.074 & \textbf{18.71} & \textbf{1.00$\times$} \\
manual     & fp16 & 2.967 & 2.929 & 3.118 & 0.120 & 19.95 & 0.85$\times$ \\
manual     & bf16 & 2.925 & 2.906 & 2.954 & 0.103 & 19.95 & 0.86$\times$ \\
sdpa       & fp32 & 2.817 & 2.814 & 2.864 & 0.039 & 19.02 & 0.89$\times$ \\
sdpa       & fp16 & 3.580 & 3.558 & 3.634 & 0.120 & 19.95 & 0.70$\times$ \\
sdpa       & bf16 & 3.622 & 3.578 & 3.794 & 0.143 & 19.95 & 0.69$\times$ \\
sdpa-flash & fp32 & 3.660 & 3.656 & 3.712 & 0.051 & 19.02 & 0.69$\times$ \\
sdpa-flash & fp16 & 4.534 & 4.522 & 4.580 & 0.109 & 19.95 & 0.55$\times$ \\
sdpa-flash & bf16 & 4.561 & 4.541 & 4.599 & 0.115 & 19.95 & 0.55$\times$ \\
\bottomrule
\end{tabular}
\caption{Inference latency and peak VRAM for \texttt{graphormer\_slim}.
Baseline is \texttt{manual / fp32}; values $<\!1\times$ are slower than
the baseline.}
\label{tab:slim}
\end{table}

%-----------------------------------------------------------------------------
\section{Discussion}

\subsection{Optimisation pay-off scales with model width, not depth}

Both architectures use 12 encoder layers, but the per-layer attention
matrix is $H \times T \times T$ with $T = N{+}1 = 33$. The
\emph{compute} delta between \texttt{base} and \texttt{slim} is
dominated by the per-layer linear projections ($d{=}768$ vs $d{=}80$,
i.e.\ $\sim$92$\times$ FLOPs in $Q,K,V$ alone), not by the attention
softmax (which scales with $T^2 H$). Fused / mem-efficient kernels save
intermediate tensor materialisation, so their absolute saving in ms
grows with kernel work. On the slim model, kernel work per call is so
small that PyTorch dispatcher overhead and the additional kernel
launches inside SDPA outweigh the saving --- producing a $0.55\times$
\emph{slowdown} for \texttt{sdpa-flash/fp16}, which is the most
optimised configuration measured.

This is a useful negative result: practitioners deploying the published
\texttt{graphormer\_slim} on edge / CPU-class devices should keep the
naive eager path; SDPA only starts paying for itself once $d \gtrsim
\mathcal{O}(10^2)$ at this $T$.

\subsection{fp16 and bf16 are within noise}

For both architectures, fp16 and bf16 are interchangeable at the
mean-latency level (within $\sim$1\% on \texttt{base}, within $\sim$2\%
on \texttt{slim}). Ampere tensor cores compute both formats at the same
peak throughput; the choice is therefore numerical, and bf16 is the
safer default for graphormer-style activations because the spatial-pos
bias can dominate $QK^{\top}$ early in training and overflow fp16 more
easily. Inference-time we recommend bf16.

\subsection{Why does autocast \emph{increase} peak VRAM?}

The peak-memory column shows fp16/bf16 at 292.5~MiB for
\texttt{base}, vs.\ 198.8~MiB for fp32 (a $\sim$47\% \emph{increase}).
This is not a measurement error. \texttt{torch.autocast} keeps weights
in fp32 and casts each operand to the low-precision dtype on the fly,
so during the cast both the fp32 source and the fp16 destination are
simultaneously resident. Genuine memory savings require calling
\texttt{model.half()} (weight-cast) rather than autocast; that path
also doubles batch capacity but is incompatible with the (small) fp32
LayerNorms inside Graphormer unless those are explicitly kept in fp32.
A weight-cast version is identified as future work.

\subsection{FlashAttention-2 vs.\ memory-efficient kernel}

On \texttt{base}, \texttt{sdpa} (auto-dispatch) and \texttt{sdpa-flash}
(forced FA-2) are within 3\% of each other. Two reasons. First,
$T{=}33$ is short --- the asymptotic advantage of FA-2 over the
memory-efficient kernel grows with $T$, and at $T{<}64$ both kernels
spend most of their time on register-tile launch and bias handling
rather than block-streamed reduction. Second, the additive bias $B$ is
large ($H \times T \times T = 32 \cdot 33 \cdot 33 = 34{,}848$
elements/sample); FA-2 supports additive bias only on Hopper
in PyTorch 2.5, so on our Ampere card the
\texttt{sdpa-flash} path silently falls back to the memory-efficient
kernel. We expect the gap to widen on H100 / B200 hardware where FA-2
admits the bias natively.

\subsection{Limitations}

(i) The 4~GB VRAM ceiling restricts \texttt{base} to $B{=}2,N{=}32$; a
full $T$-scan up to $N{=}128$ (the model's
\texttt{max\_nodes}) was not feasible. (ii) We do not measure with
\texttt{torch.compile} --- adding it is a single-line change and
should compose multiplicatively with SDPA. (iii) Synthetic batches
exercise the full computational graph but assume a uniformly-distributed
spatial-pos matrix; on real ZINC molecules the bias is sparser (most
distances are small integers) which may favour the math kernel slightly.
(iv) We did not load the official PCQM4Mv1/v2 checkpoint into the
end-to-end timed loop --- the loader is implemented and tested for
key-mapping correctness, but adds I/O cost orthogonal to kernel
choice, so we keep it out of the ablation.

%-----------------------------------------------------------------------------
\section{Conclusion}

We extracted the Graphormer encoder and structural-encoding stack into a
dependency-free PyTorch~2 module, integrated it with the official
Microsoft checkpoints via a state-dict re-mapping utility, and
benchmarked three attention back-ends across three precisions on an
Ampere GPU. The headline result is that
\texttt{sdpa+fp16}/\texttt{bf16} delivers a $2.09\times$ speed-up on
\texttt{graphormer\_base} at no accuracy cost (the bias is preserved
bit-for-bit), while the same recipe is a slowdown on
\texttt{graphormer\_slim} because dispatch overhead exceeds the
arithmetic saved. KV-caching is excluded with a structural argument:
Graphormer is a single-shot bidirectional encoder whose attention bias
does not factorise along any sequential axis. Promising next steps are
(i) weight-cast inference for genuine VRAM reduction, (ii)
\texttt{torch.compile} + CUDA Graphs on top of SDPA, and (iii)
re-running the same matrix on H100 to quantify the FA-2 bias path that
is unavailable on Ampere.

%-----------------------------------------------------------------------------
\bibliographystyle{plain}
\bibliography{references}

\end{document}
